\newpage
\section{Hypothesis: Topology-Aware Operation as a Precondition for Data-Safe Scale-In}
\label{sec:hypothesis}

The limitations of topology-blind, count-based scaling described in
Section~\ref{subsec:stateful-challenges} motivate the central hypothesis of this thesis:
that \textit{topology-aware operation}---the encoding of Redis-specific resharding logic
into a Kubernetes operator---is a precondition for the data-safe scale-in of a sharded,
stateful workload, and that the native StatefulSet contraction performed by the Horizontal
Pod Autoscaler, being blind to the cluster's slot map, cannot preserve data when a
slot-owning master is removed. The hypothesis is therefore framed not in terms of how
\textit{quickly} or how \textit{autonomously} a scaling decision is reached, but in terms of
whether the scaling \textit{action}, once taken, leaves the cluster's keyspace intact; the
governing metric is accordingly \textit{data loss}---keys lost and hash slots left
permanently uncovered---rather than responsiveness or resource efficiency.

\subsection{The Operator Pattern as a Foundation for Stateful Automation}

The Kubernetes \textit{operator pattern} was conceived precisely to bridge the gap between
the generic automation that Kubernetes provides out of the box and the deep,
application-specific knowledge required to manage complex stateful systems reliably.
According to the official Kubernetes documentation, the operator pattern \enquote{aims to
capture the key aim of a human operator who is managing a service or set of services},
encoding that operational expertise into a software controller that participates in the
standard reconciliation loop (\cite{kubernetes:operatorpattern}). Operators extend the
Kubernetes API through \textit{Custom Resource Definitions} (CRDs) together with an
accompanying controller that continuously reconciles the observed state of the custom
resource against its declared desired state.

Critically, the CNCF Operator White Paper explicitly acknowledges that \enquote{Kubernetes
primitives were not built to manage state by default} and that \enquote{relying on
Kubernetes primitives alone brings difficulty managing stateful application} lifecycle
(\cite{cncf:operatorwhitepaper}). The operator pattern is therefore not merely a convenience
but the architecturally recommended means of encoding domain-specific invariants---data
resharding, replication-topology maintenance, and quorum-safe scaling---into a
Kubernetes-native automation layer (\cite{k8s:operator}). For a Redis Cluster in particular,
a custom operator is able to express and enforce the one invariant that a generic,
count-based controller fundamentally cannot: that a departing master must first be
\textit{drained}---its hash slots migrated onto the surviving masters---and only then
removed (\cite{opstree:redisoperator}; \cite{redis:scaling}). It is this single invariant,
rather than any property of how the scale-in happens to be triggered, on which the
data-safety of the operation turns.

\subsection{The Scaling Mechanism, Not the Trigger, Governs Data Safety}

It is useful to separate two concerns that discussions of autoscaling frequently conflate.
The first is the \textit{trigger}: the question of when, and on the basis of which signal, a
scaling decision is taken---the axis along which the reactive HPA, which acts on lagging CPU
and memory utilisation, is contrasted with proactive, event-driven approaches such as KEDA
that act on application-level metrics ahead of resource exhaustion
(\cite{keda:scalingdeployments}). The second is the \textit{mechanism}: the question of how
the resulting change in size is actually carried out against the running cluster. For a
sharded, stateful store these two concerns are orthogonal, and the present thesis contends
that, for scale-in, it is the mechanism alone that determines whether the data survives.

The argument is straightforward. However intelligently the timing of a scale-in is chosen,
if the removal is executed by deleting a pod by ordinal with no accompanying
reshard---which is precisely what a StatefulSet contraction does---then the departing
master's slots are abandoned and the keys hashing to them are lost; the quality of the
trigger cannot rescue an unsafe mechanism. Conversely, a topology-aware drain preserves the
keyspace irrespective of whether it was set in motion by a reactive threshold, a proactive
metric, or a manual operator action. Because the data-safety question is in this way decided
by the mechanism and not by the trigger, this thesis deliberately holds the trigger
constant---indeed, removes it altogether, driving the size change by hand---in order to
isolate the mechanism as the sole variable. The event-driven trigger layer, and KEDA in
particular, is on this reading an enhancement to \textit{when} scaling occurs that is
orthogonal to \textit{whether} it is safe, and it is accordingly reserved for future work
rather than examined here.

\subsection{Positioning Within the Literature}

The existing body of autoscaling research has concentrated overwhelmingly on the trigger
problem---on making scaling decisions better-timed and more anticipatory. Mondal et al.\
identify the HPA's reliance on CPU and memory as lagging indicators that render it
\enquote{incapable of foreseeing upcoming workload spikes}, leading to Quality-of-Service
violations, long-tail latency, and wasted resources, while the NimbusGuard study of
Wanigasooriya and Ekanayake demonstrates that proactive autoscaling frameworks
\enquote{translate into superior performance and cost efficiency compared to existing
reactive methods} under identical, phased load patterns; Toka et al.\ similarly enumerate
the structural drawbacks of purely observation-based, manually tuned reactive scaling
(\cite{arxiv:nimbusguard}). This literature establishes, convincingly, the value of
improving \textit{when} a stateful workload is scaled.

What it largely leaves unexamined is whether the scaling \textit{action} itself is safe for a
sharded store---the precise concern of the present work. Here the relevant grounding is
provided not by the autoscaling-performance literature but by the operator and
Redis-cluster sources: the CNCF Operator White Paper's observation that Kubernetes
primitives were not built to manage state (\cite{cncf:operatorwhitepaper}), the Redis
Cluster specification's requirement that the \num{16384} hash slots be migrated rather than
merely abandoned when a node departs (\cite{redis:clusterspec}), and the recognised risk of
performing that resharding incorrectly (\cite{googlecloud:zerodowntime}). That the
automation of Redis lifecycle management by operators remains only partial, and an open
challenge, has itself been noted by the CNCF community (\cite{cncf:redisoperator}). The
hypothesis advanced here therefore occupies a gap that the timing-focused literature leaves
open: it asks not how to scale a stateful workload more cleverly, but whether a given
scaling mechanism preserves its data at all.

\subsection{Scope and Expected Outcome}

Given this foundation, the thesis hypothesises a sharp, qualitative divergence between the
two mechanisms when a single slot-owning master is removed from an otherwise healthy
cluster. The native StatefulSet contraction---the action the HPA performs on
scale-down---is expected to produce a structural \textit{data cliff}: with neither a drain
nor a reshard, the departing master's slots are left uncovered, the cluster transitions to a
failed state with fewer than \num{16384} slots covered, and a measurable fraction of the
keys is irrecoverably lost. Crucially, this loss is expected to occur even under zero client
load, which would demonstrate that it is structural rather than a consequence of write
pressure or unfortunate timing. The operator-driven removal, by contrast, is expected to
drain the departing master gracefully, holding continuous \num{16384}-slot coverage and
preserving every key.

This expectation is, however, deliberately qualified. The operator is hypothesised to be
\textit{necessary but not, on its own, sufficient} for fully hands-off elastic operation: its
absence of self-healing for an interrupted reshard is expected to surface, in some runs, as a
recoverable stall in which a slot is left mid-migration and a human operator must intervene
with \texttt{redis-cli -{}-cluster fix} to settle the cluster. The decisive point is that such
a stall, while a genuine operational cost, is categorically distinct from the native
mechanism's data cliff, because no data is lost. The experimental evaluation presented in
Chapter~\ref{sec:results} is designed to test these expectations directly, by performing a
controlled, quiescent, single-master scale-in against each mechanism and measuring the
outcome through the cluster's own slot and key telemetry, with data loss as the decisive
criterion.
